\renewcommand{\Anno}{2008}

\renewcommand{\NumeroProblema}{3}

\renewcommand{\Testo}{ 
	Una colonna dell'Enalotto \`e una sequenza
	di $12$ simboli, scelti tra \texttt{1}, \texttt{X}, \texttt{2}.
	\begin{enumerate}
		\item Quante sono le colonne dell'Enalotto che presentano una
		ed una sola volta due simboli consecutivi uguali?

		\item Quante sono le differenti colonne dell'Enalotto che
		cominciano e finiscono con $1$, e che non presentano mai due
		simboli consecutivi uguali? 
	\end{enumerate}
}

\renewcommand{\Risposta}{
\begin{enumerate}
	\renewcommand{\labelenumi}{(\alph{enumi})} 
	\item  33792

	\item  682
\end{enumerate}
}

\renewcommand{\Aiutino}{
\begin{enumerate}
\renewcommand{\labelenumi}{(\alph{enumi})} 
	\item Trovare quante sono le colonne con 11 simboli che
	non presentano mai 2 simboli consecutivi uguali: in quanti modi
	possiamo scegliere il primo simbolo? In quanti modi possiamo
	scegliere ciascuno dei simboli successivi?
	
	A questo punto non resta che scegliere uno degli 11 simboli e
	raddoppiarlo. In quanti modi possiamo farlo?

	\item Indichiamo con $U_{n}$ il numero di colonne che iniziano e
	terminano per \texttt{1} e non hanno due simboli uguali
	consecutivi.  Indichiamo con $D_{n}$ il numero di colonne che
	iniziano per \texttt{1}, terminano per \texttt{X} o \texttt{2}, e
	non hanno due simboli uguali consecutivi.  Segnaliamo 2 possibili
	modi di procedere.
	\begin{itemize}
		\item  	Mostrare che valgono le relazioni ricorrenti
		$$U_{n+1}=D_{n},
		\hspace{3em}
		D_{n+1}=D_{n}+2U_{n}.$$
		
		Calcolati $U_{3}$ e $D_{3}$ si pu\`{o} ripetere la doppia
		iterazione fino ad arrivare a $U_{12}$.
		
		\item  Mostrare che $U_{n}+D_{n}=2^{n-1}$ e quindi
		$U_{n+1}=D_{n}=2^{n-1}-U_{n}$. Da questa si deduce che
		$$U_{12}=2^{10}-2^{9}+2^{8}-2^{7}+\ldots-2.$$
	\end{itemize}
	
	Entrambi i procedimenti permettono anche di arrivare alla formula
	generale per $U_{n}$, che risulta essere 
	$$U_{n}=\frac{2^{n-1}+2\cdot(-1)^{n-1}}{3}.$$

\end{enumerate}
}

\renewcommand{\SoluzioneUfficiale}{\DomandaA Le colonne richieste sono
$11\cdot 3\cdot 2^{10}=33792$.  Infatti una colonna che contiene una
ed una sola volta due simboli consecutivi uguali si ottiene in modo
unico prendendo una successione di 11 simboli, che non ha mai due
simboli consecutivi uguali, e raddoppiando uno dei simboli.  Si hanno
evidentemente 11 possibili scelte per il simbolo da raddoppiare,
dunque basta trovare il numero delle successioni di lunghezza 11 che
non hanno mai due simboli consecutivi uguali.  Si hanno ovviamente 3
scelte per il primo simbolo, e 2 scelte per ognuno dei simboli
successivi, per un totale di $3\cdot 2^{10}$ possibilit\`{a}.

\DomandaB Indichiamo con $U_{n}$ il numero di colonne che iniziano e
terminano per \texttt{1} e non hanno due simboli uguali consecutivi.
Indichiamo con $D_{n}$ il numero di colonne che iniziano per
\texttt{1}, terminano per \texttt{X} o \texttt{2}, e non hanno due
simboli uguali consecutivi.

Mostriamo che per ogni intero $n\geq 2$ si ha che
\begin{eqnarray}
    D_{n+1} & = & D_{n}+2U_{n},
    \label{eq:ric-dn}  \\
    U_{n+1} & = & D_{n}.
    \label{eq:ric-un}
\end{eqnarray}

Per dimostrare tali formule iniziamo ad indicare con $\mathcal{D}_{n}$
l'insieme delle sequenze che vengono conteggiate da $D_{n}$, e con
$\mathcal{U}_{n}$ l'insieme delle sequenze che vengono conteggiate
da $U_{n}$. Un elemento di $\mathcal{D}_{n+1}$ pu\`{o} essere ottenuto
in due modi:
\begin{itemize}
    \item  prendendo un elemento di $\mathcal{U}_{n}$ e aggiungendo un
    simbolo a scelta tra \texttt{X} e \texttt{2};

    \item  prendendo un elemento di $\mathcal{D}_{n}$ e aggiungendo
    l'unico simbolo possibile (non possiamo n\'{e} mettere \texttt{1},
    n\'{e} ripetere l'ultimo simbolo).
\end{itemize}

In questo modo si ottengono tutti gli elementi di $\mathcal{D}_{n+1}$,
e quelli prodotti nel primo modo sono diversi da quelli prodotti nel
secondo modo. Questo giustifica la (\ref{eq:ric-dn}).

Un elemento di $\mathcal{U}_{n+1}$ pu\`{o} invece essere ottenuto in
un solo modo:
\begin{itemize}
    \item prendendo un elemento di $\mathcal{D}_{n}$ e aggiungendo
    \texttt{1} in fondo.
\end{itemize}

Questo giustifica la (\ref{eq:ric-un}). 

\`{E} ora evidente che gli elementi di $\mathcal{U}_{3}$ sono
\texttt{1X1} e \texttt{121}, mentre quelli di $\mathcal{D}_{3}$ sono
\texttt{12X} e \texttt{1X2}, e pertanto $U_{3}=D_{3}=2$.  Applicando
con pazienza le relazioni ricorrenti troviamo pian piano i valori di
$U_{n}$ e di $D_{n}$ indicati nella seguente tabella.
\begin{center}
	\renewcommand{\arraystretch}{1.2}
	\begin{tabular}{|c||c|c|c|c|c|c|c|c|c|c|}
		\hline
		$n$ & 3 & 4 & 5 & 6 & 7 & 8 & 9 & 10 & 11 & 12  \\
		\hline\hline
		$U_{n}$ & 2 & 2 & 6 & 10 & 22 & 42 & 86 & 170 & 342 & 682  \\
		\hline
		$D_{n}$ & 2 & 6 & 10 & 22 & 42 & 86 & 170 & 342 & 682 &   \\
		\hline
	\end{tabular}
\end{center}

In conclusione $U_{12}=682$.
}

\renewcommand{\SoluzioniAlternative}{\TitoloSoluzioneAlternativaA
\DomandaB 

Utilizziamo le stesse notazioni della soluzione precedente.
Osserviamo preliminarmente che $U_{n}+D_{n}=2^{n-1}$, in quanto
$\mathcal{U}_{n}\cup\mathcal{D}_{n}$ \`{e} l'insieme di tutte le
colonne di $n$ simboli che iniziano con \texttt{1} e non hanno mai 2
simboli consecutivi (e queste sono appunto $2^{n-1}$ poich\'{e} ci
sono 2 possibilit\`{a} per tutti i simboli successivi al primo).
Inoltre $U_{n+1}=D_{n}$, cos\`{i} come dimostrato nella soluzione
precedente.

Mettendo insieme le 2 relazioni troviamo che 
\begin{equation}
	U_{n+1}=2^{n-1}-U_{n}, 
	\label{eq:ric-somma}
\end{equation}
e pertanto
$$U_{12}=2^{10}-U_{11}=2^{10}-2^{9}+U_{10}=
2^{10}-2^{9}+2^{8}-U_{9}=\ldots$$

Procedendo in questo modo fino ad arrivare ad $U_{3}$, che fa 2, si ha
la formula
$$U_{12}=2^{10}-2^{9}+2^{8}-2^{7}+\ldots+2^{2}-2,$$
che si calcola facilmente.
% 
% \TitoloSoluzioneAlternativaB \DomandaB Spezziamo una colonna da 12 in
% 2 semicolonne da 6.  Indichiamo con $S_{1}$, $S_{x}$, $S_{2}$ il
% numero di (semi)colonne, composte da 6 simboli, che iniziano con
% \texttt{1}, non contengono mai 2 simboli consecutivi uguali, e
% finiscono con \texttt{1}, \texttt{X}, \texttt{2}, rispettivamente.
% 
% Allora il numero richiesto \`{e}
% $$2(S_{1}S_{x}+S_{x}S_{2}+S_{2}S_{1}),$$
% in quanto ognuno dei tre tipi di semicolonna pu`o essere preso come
% parte iniziale o come parte finale (rovesciata).
% 
% Per calcolare $S_{1}$ si pu\`{o} ripetere il ragionamento precedente,
% cio\`{e} spezzare la colonna da 6 in due semicolonne da 3.  Poich\'{e}
% le possibili semicolonne di lunghezza 3 sono soltanto 4 (\texttt{1X2},
% \texttt{1X1}, \texttt{121}, \texttt{12X}) si ottiene facilmente che
% $$S_{1}=2(2\cdot 1+1\cdot 1+1\cdot 2)=10.$$
% 
% Per calcolare $S_{x}$ ed $S_{2}$ osserviamo intanto che
% 
% 
% 
% 
% Per calcolare i valori di s1, sx, s2 si osservi anzitutto che per s1
% si pu`o ripetere il ragionamento precedente, tenendo conto che le
% semicolonne di lunghezza 3 che cominciano per 1 e non presentano mai
% due simboli consecutivi uguali sono soltanto 
}

\renewcommand{\Osservazioni}{\TitoloOsservazione 

Questo problema offre parecchi spunti.
\begin{enumerate}
	\item Manipolando algebricamente la (\ref{eq:ric-un}) e la
	(\ref{eq:ric-dn}) otteniamo che
	\begin{equation}
		U_{n+2}=D_{n+1}=D_{n}+2U_{n}=U_{n+1}+2U_{n}.
		\label{eq:ric2-un}
	\end{equation}
	
	Questa \`{e} una ricorrenza che coinvolge una sola delle 2
	successioni.  Dalla teoria generale delle successioni per
	ricorrenza lineari segue immediatamente che $U_{n}=\alpha
	2^{n}+\beta(-1)^{n}$ (dove $2$ e $-1$ sono le soluzioni di
	$x^{2}=x+2$).  A questo punto $\alpha$ e $\beta$ si ricavano
	imponendo che $U_{2}=0$ e $U_{3}=2$.  Si ottiene cos\`{i} la
	formula generale
	\begin{equation}
		U_{n}=\frac{2^{n-1}+2\cdot(-1)^{n-1}}{3}.
		\label{eq:un}
	\end{equation}
	
	\item Dalla (\ref{eq:ric-somma}) si ottiene, con una facile
	induzione, la formula
	$$U_{n}=2^{n-2}-2^{n-3}+2^{n-4}-\ldots+(-1)^{n+1}\cdot 2,$$
	e questa si calcola esplicitamente come somma di una progressione
	geometrica di ragione $-2$. Ovviamente in questo modo si riottiene
	la (\ref{eq:un}).

	\item La formula (\ref{eq:ric2-un}) ammette, oltre alla
	giustificazione algebrica presentata, anche una giustificazione
	diretta di tipo combinatorio.  Infatti un elemento di
	$\mathcal{U}_{n+2}$ pu\`{o} essere ottenuto in due modi (che si
	escludono a vicenda):
	\begin{itemize}
		\item prendendo un elemento di $\mathcal{U}_{n+1}$, sostituendo
		l'\texttt{1} finale con l'unico altro simbolo possibile, e
		aggiungendo un \texttt{1} alla fine;

		\item prendendo un elemento di $\mathcal{U}_{n}$ e aggiungendo
		prima un simbolo a scelta tra \texttt{X} e \texttt{2}, e poi
		un \texttt{1}.
	\end{itemize}
\end{enumerate}
}

\renewcommand{\Commenti}{
}